\section{Experimental Costs for Key Recovery}
\label{sec:total-costs}

Experiments were run on a single-socket, single-NUMA-node x86\_64 machine with
an Intel Core Ultra 7 265H processor (family 6, model 197, stepping 2): 16
cores and 16 hardware threads (one thread per core), with a maximum frequency
of 5.3\,GHz, 30\,GiB of RAM and 8\,GiB of swap, a 24\,MiB shared L3 cache,
28\,MiB of L2 cache in total, and 480\,KiB of L1 data cache and 768\,KiB of L1
instruction cache in total.

All experiments below were successful in key recovery. Code will be available on Github shortly if the 
author gets around to it. 

\subsection{Parameter Sets Tested}\label{sec:param-sets}

In addition to the ID-0 (128-bit) parameter set from the RLCE specification, we
also tested several smaller parameter sets, and present them and name them here
to avoid cluttering the timing table below; all sets below have $w=n-k$. 


\begin{table}[!ht]
\caption{Parameter Sets Tested}
\begin{center}
\begin{tabular}{llcccc}
\toprule
Name        & q & n & k & $n'$\\
\midrule
256  & 256 & 198 & 165 & 90\\
512  & 512 & 500 & 375 & 293\\
\bottomrule
\end{tabular}
\end{center}
\end{table}

\subsection{Timing tables}\label{sec:timing}

Below, phase 1 refers to the total time taken summed across threads.
\begin{table}[!ht]
\caption{Timing Results}
\begin{center}
    \footnotesize
    \begin{tabular}{llrrrrr}
        \toprule
        Set                           & Trials             & Threads & Phase~1 Time & Phase~2 Time \\
        \midrule
        256 & 50 & 10 & 288.27 (s) & 0.02 (s)\\
        512 & 44 & 12 & 37278.15 (s) & 0.09 (s) \\
        ID-0    & 41             & 12   & 109673.21 (s)        & 0.23 (s)     \\
        \bottomrule
    \end{tabular}
\end{center}
\end{table}